%%%%%%%%%%%%Outline%%%%%%%%%%%%%%%%%%

%%% Section 5 \lang: Hardware-Idiom Operations

%%%
% message<\lang is the seven-operation instance of \tool for hardware idioms>
% with \tool we build \lang, seven operations for hardware idioms
% each operation targets an idiom with no single-instruction eBPF form
% the rotate, for example, covers its 64-bit and 32-bit variants through descriptors
% the section closes with the Lean 4 proofs that supply the owed evidence

%%%
% message<the seven operations are a deliberate design point, with payoffs>
% the seven operations are a deliberate choice within the space \tool leaves open
% the idioms of the characterized kind are few
% a handful of modules therefore covers the kernel JIT's clearest losses
% each proof sequence stays a small fragment of the program
% the equivalence evidence each operation owes covers only that fragment
% the compiler spells six of the seven out as fixed patterns the recognizer matches in place
% the prefetch alone has no eBPF form of any length
% the recognizer instead inserts the prefetch at selected sites, with a no-op proof sequence

%%%
% message<the table lists the operations; native emits differ in shape and availability>
% the table lists the seven operations with proof sequence and native emit per architecture
% the native emits differ across architectures, in shape and in availability
% a conditional select is one extended instruction on x86-64, two on ARM64 (TST, CSEL)
% the load-fused byte swap splits the same way on ARM64, into LDR and REV
% one operation, the address computation LEA, exists on x86-64 alone

%%%
% message<we prove in Lean 4 that each native emit equals its proof sequence>
% the dual lowering imposes a proof obligation
% every safety property on the eBPF expansion must transfer to the native sequence
% we prove in Lean 4 the two compute the same function on the destination register
% both sides are pure computations, so functional equivalence
% each ISA gets a small-step semantics, a program is the fold of transitions
% the native instruction's semantics is ported from an existing formalization
% the cross-ISA proof is factored through a hardware-independent specification
% per-lang bitvector spec, ISA-local lemmas refine the spec on the destination register

%%%%%%%%%%%%Outline%%%%%%%%%%%%%%%%%%

\section{\lang: Hardware-Idiom Operations}
\label{sec:lang}
% \lang：硬件习语操作

With \tool, we build \lang, a set of seven operations for hardware idioms.
% 借助 \tool，我们构建了 \lang，一组七个硬件习语操作。
Each operation targets an idiom for which the eBPF instruction set has no single-instruction form.
% 每个操作针对 eBPF 指令集没有单指令形式的习语。
The rotate, for example, covers its 64-bit and 32-bit variants through separate descriptors.
% 例如旋转通过单独的描述符覆盖其 64 位和 32 位变体。
The section closes with the Lean 4 proofs that supply the evidence \S\ref{subsec:safety} requires.
% 本节以提供所需证据的 Lean 4 证明结束。

The seven operations are a deliberate choice within the space that \S\ref{subsec:adding} leaves open.
% 这七个操作是在所留空间内的刻意选择。
The idioms of the kind \S\ref{subsec:char-idiom} characterizes are few.
% 所特征化的这类习语很少。
A handful of modules therefore covers the kernel JIT’s clearest losses.
% 因此少数模块就覆盖了内核 JIT 最明显的损失。
Each proof sequence stays a small fragment of the program.
% 每个证明序列都是程序的一小片段。
The equivalence evidence each operation owes covers only that fragment.
% 每个操作所需的等价证据只覆盖该片段。
The compiler emits six of the seven idioms as fixed bytecode patterns, which the recognizer matches in place.
% 编译器将七个习语中的六个发射为固定的字节码模式，识别器就地匹配。
The prefetch alone has no eBPF form of any length.
% 只有预取没有任何长度的 eBPF 形式。
The recognizer instead inserts the prefetch at sites it selects, with a no-op proof sequence.
% 识别器在其选择的位置插入预取，使用空操作证明序列。

Table~\ref{tab:descriptors} lists the seven operations, each with its proof sequence and its native emit on x86-64 and ARM64.
% 表列出了七个操作，每个都有其证明序列及其在 x86-64 和 ARM64 上的本地发射。
The native emits differ across architectures, in shape and in availability.
% 本地发射在架构间的形状和可用性上有所不同。
A conditional select compiles to one extended instruction on x86-64, emitting the \texttt{CMP+CMOV} pair, and to two extended instructions on ARM64, one emitting \texttt{TST} and one emitting \texttt{CSEL}.
% 条件选择在 x86-64 上编译为一条扩展指令，发射 CMP+CMOV 对，在 ARM64 上编译为两条扩展指令，一条发射 TST，一条发射 CSEL。
The byte swap, fused with its load, splits the same way on ARM64, into an \texttt{LDR} and a \texttt{REV}.
% 与加载融合的字节交换在 ARM64 上以相同方式分割，成为 LDR 和 REV。
One operation, the address computation \texttt{LEA}, exists on x86-64 alone.
% 一个操作，地址计算 LEA，仅存在于 x86-64。

\begin{table}[t]
\centering
\small
\caption{The seven \lang operations, each with its typical proof sequence and the native code per architecture. Exact proof sequences vary by operation variant and architecture (\S\ref{subsec:pipeline}). A dash marks an architecture with no native code, where the operation falls back to its proof sequence. Each native-code cell lists the instructions emitted for one operation instance. Table~\ref{tab:isa-gap} counts the instructions along each compilation path for the idioms the two tables share, and the native compiler's instruction choice there can differ from the native code here.}
\label{tab:descriptors}
\begin{tabular}{@{}llll@{}}
\toprule
 & & \multicolumn{2}{c}{Native code} \\
\cmidrule(lr){3-4}
Operation & Proof sequence & x86-64 & ARM64 \\
\midrule
Rotate & $n$ one-bit rotations & \texttt{ROL} & \texttt{ROR} \\
Conditional select & branch + move & \texttt{CMP+CMOV} & \texttt{TST+CSEL} \\
Bit extract & shift + and & \texttt{BEXTR} & \texttt{UBFX} \\
Byte swap & load + byte swap & \texttt{MOVBE} & \texttt{LDR+REV} \\
Prefetch & no-op (\texttt{JA +0}) & \texttt{PREFETCHT0} & \texttt{PRFM} \\
Bulk copy & loads + stores & \texttt{MOV} & \texttt{LDP+STP} \\
LEA & shift + adds & \texttt{LEA} & -- \\
\bottomrule
\end{tabular}
\end{table}

\noindent\textbf{Formal verification.}\label{sec:formal-verification}
% \noindent\textbf{形式化验证。}
The dual lowering of a \lang: one path expands to vanilla eBPF for the verifier, the other emits native code in the JIT, imposes a proof obligation.
% \lang 的双重降低：一条路径为验证器扩展为原生 eBPF，另一条在 JIT 中发射本地代码，这施加了证明义务。
The native sequence that the CPU actually executes must produce the same eBPF-observable architectural state as the eBPF expansion.
% CPU 实际执行的本地序列必须产生与 eBPF 扩展相同的 eBPF 可观察架构状态。
We therefore prove, in Lean 4, that for each supported \lang operation, the eBPF expansion and the corresponding native sequence produce the same eBPF-observable architectural state.
% 因此我们在 Lean 4 中证明，对于每个支持的 \lang 操作，eBPF 扩展和相应的本地序列产生相同的 eBPF 可观察架构状态。
For register-only operations such as rotate, conditional select, and bit extract, the refinement reduces to equality on the destination register.
% 对于仅寄存器操作如旋转、条件选择和位提取，精化归结为目标寄存器上的相等。
For memory-touching operations such as byte swap (fused with its load) and bulk copy, the proof obligation is equality on the projected architectural state over live registers and the accessed memory, leaving all other eBPF-visible memory unchanged.
% 对于涉及内存的操作如字节交换（与加载融合）和批量复制，证明义务是活寄存器和访问内存上投影架构状态的相等，其他 eBPF 可见内存保持不变。
Prefetch is modeled as an architectural no-op; microarchitectural effects such as cache state are outside the semantic model.
% 预取被建模为架构空操作；缓存状态等微架构效果在语义模型之外。
We give each ISA a small-step register-transfer semantics; a program denotes the fold of its per-instruction transitions.
% 我们给每个 ISA 一个小步寄存器传递语义；程序表示其逐指令转换的折叠。
The native instruction’s semantics is ported from the existing formalization~\cite{lnsym}.
% 本地指令的语义从现有形式化移植。
The cross-ISA proof is factored through an intermediate hardware-independent specification.
% 跨 ISA 证明通过中间的硬件无关规范分解。
The refinement is stated for verifier-accepted proof sequences and verifier-safe initial states.
% 精化针对验证器接受的证明序列和验证器安全的初始状态陈述。
For each \lang operation, we define a specification over the relevant architectural state and discharge ISA-local correctness lemmas stating that the eBPF expansion and the native sequence each refine the specification.
% 对于每个 \lang 操作，我们定义相关架构状态上的规范，并卸载声明 eBPF 扩展和本地序列各自精化规范的 ISA 局部正确性引理。
% This keeps the per-ISA obligations local, and we therefore accept a new \lang by supplying its specification and refinement proofs.
